% Generated draft: requires mathematical and source review.
\documentclass[11pt]{article}
\usepackage[margin=25mm]{geometry}
\usepackage{fontspec}
\setmainfont{Latin Modern Roman}
\usepackage{amsmath,amssymb,mathtools,unicode-math}
\setmathfont{Latin Modern Math}
\usepackage{enumitem,needspace,xurl,hyperref}
\setlength{\parindent}{0pt}
\setlength{\parskip}{5pt}
\setlength{\emergencystretch}{3em}
\newcommand{\N}{\mathbb{N}}
\newcommand{\Z}{\mathbb{Z}}
\newcommand{\Q}{\mathbb{Q}}
\newcommand{\R}{\mathbb{R}}
\newcommand{\C}{\mathbb{C}}
\newcommand{\F}{\mathbb{F}}
\newcommand{\A}{\mathbb{A}}
\newcommand{\PP}{\mathbb{P}}
\newcommand{\E}{\mathbb{E}}
\newcommand{\eps}{\varepsilon}
\newcommand{\dd}{\,\mathrm{d}}
\begin{document}
\section*{UnsolvedMath ID: 700940}
\textbf{Source problem code:} catalogue entry
\subsection{Definitions and mathematical statement}
A group is \emph{polycyclic} if it has a finite subnormal series whose factors are cyclic, with finite cyclic factors allowed. A group is \emph{locally nilpotent} if every finitely generated subgroup is nilpotent. For subgroups \(A,B\leq G\), the equality \(G=AB\) means that every \(g\in G\) can be written as \(g=ab\) with \(a\in A\) and \(b\in B\); neither subgroup is assumed normal.

The source calls \(G\) \emph{radical} when it has an ascending normal series with locally nilpotent factors. Under the continuous, possibly transfinite interpretation of a normal series, this means that there are an ordinal \(\lambda\) and subgroups \((G_\alpha)_{\alpha\leq\lambda}\) such that
\[
G_0=1,\qquad G_\lambda=G,\qquad G_\alpha\trianglelefteq G,
\]
\[
G_\alpha\leq G_\beta\quad(\alpha\leq\beta\leq\lambda),\qquad
G_\delta=\bigcup_{\alpha<\delta}G_\alpha
\quad\text{for every nonzero limit ordinal }\delta\leq\lambda,
\]
and \(G_{\alpha+1}/G_\alpha\) is locally nilpotent for every \(\alpha<\lambda\). The supplied excerpt does not explicitly state the series conventions; the displayed interpretation, including normality in the whole group, requires confirmation against the source.

\textbf{Question (F. de Giovanni, 1995).} For every group \(G\) and every pair of polycyclic subgroups \(A,B\leq G\), if \(G=AB\) and \(G\) has an ascending normal series with locally nilpotent factors, must \(G\) be polycyclic? No finiteness assumption on \(G\), beyond what follows from the stated hypotheses, is imposed.
\subsection{Short English statement}
F. de Giovanni asks whether a radical group that is the setwise product of two polycyclic subgroups must itself be polycyclic. The precise convention for the ascending normal series should be checked against the original source.
\subsection{Sources}
\begin{itemize}
\item Local UnsolvedMath registry, permanent ID \(700940\), exact source problem code \texttt{catalogue entry}, title \emph{Kourovka Notebook Problem 13.22}. The supplied registry labels the problem \emph{open}; this is a supplied status, not an independently verified mathematical or literature assessment. The record was updated on 2026-10-16.
\item \emph{The Kourovka Notebook (Unsolved Problems in Group Theory)}, original problem 13.22, 13th issue (1995); also identified in the supplied record with the 21st issue edition (2026), \texttt{arXiv:1401.0300v48}. Proposed by F. de Giovanni. The supplied collection description names E. I. Khukhro and V. D. Mazurov as editors.
\end{itemize}
\end{document}
